(1) 取Bloch基 $(a_k,b_k)$，可写 $H(k)=-\begin{pmatrix}0&v+we^{-ik}\\v+we^{ik}&0\end{pmatrix}$，能量为 $\boxed{E_\pm(k)=\pm\sqrt{v^2+w^2+2vw\cos k}}$。矩阵没有对角项，满足 $\sigma_zH\sigma_z=-H$ 的手征对称。

(2) 在Brillouin区 $k\in[-\pi,\pi]$，复数 $v+we^{ik}$ 描绘以正实轴 $v$ 为中心、半径 $w$ 的圆。$w>v$ 时曲线包围原点，绕数绝对值为1；$w<v$ 时不包围，绕数0。最小正带隙是 $|v-w|$，两带间隙为 $\boxed{\Delta=2|v-w|}$，仅在 $v=w$、$k=\pi$ 闭合。

(3) 半无限左端A轨道终止，在零能寻找所有B振幅为0的态。对B点的Schrödinger方程为 $v A_j+wA_{j+1}=0$，所以 $A_{j+1}=-(v/w)A_j$。因而 $\boxed{\psi_{A,j}=\sqrt{1-(v/w)^2}(-v/w)^{j-1},\quad\psi_{B,j}=0}$，仅当 $v<w$ 时归一化。局域长度为 $\xi=1/\ln(w/v)$（以晶胞为单位）。

(4) 穿越 $v=w$ 时体能隙闭合，绕数改变；开边界的拓扑相存在局域零模。有限链的左右端尾部重叠使两零模形成指数小的能量劈裂，量级随长度按 $(v/w)^N$ 下降。若增加不同子格的现场能差而破坏手征对称，零模一般不再严格钉在零能。
